\documentclass{report}
\usepackage{graphicx}
\usepackage{blindtext}
\usepackage{geometry}
\usepackage{caption}
\usepackage{subcaption}
\usepackage{geometry}
\usepackage{amsmath}
\usepackage{amssymb}
\usepackage{tabularx}
\usepackage{lipsum}
\usepackage{xcolor}
\renewcommand{\arraystretch}{1.2}

\geometry{
	a4paper,
	right=17mm,
	bottom=20mm,
	left=17mm,
	top=15mm,
}
\title{Independent Research \& Innovation for Students \\ (IRIS) }

\author{
Prof. Dr. Torben Ferber \\ Supporting Professor 
\and 
Prof. Dr. Markus Klute \\ Supporting Professor 
\and 
Prof. Dr. Frank Simon \\ Supporting Professor 
\and
Sören Finna \\ BSc Student
\and 
Vanessa Breitenstein \\ BSc Student
\and 
Marcus Lai \\ MSc Student}

\date{Lecture-free Period 2024}

\begin{document}
\begin{figure}
	\includegraphics[width=0.4\linewidth]{KIT}
	\centering
	\label{KITlogo}
\end{figure}

\maketitle
\tableofcontents
\setcounter{chapter}{-1}

\chapter{Preamble}

	\section{What is IRIS} 
        IRIS (Independent Research \& Innovation for Students) is a new initiative to allow keen BSc or MSc students to explore and self-propose pertinent, unique, and ambitious short term projects. These projects are not from research groups in KIT, but instead are to be identified and proposed by student groups themselves.
        \newline
    
        \noindent
        This opportunity presents a low barrier to entry for participants to get introduced to research projects, while allowing them to develop their own research interests. One of the aims of IRIS is as an approachable pathway that directly contrasts student research (Bachelor's thesis, Master's Thesis, HIWI Jobs) in KIT. Open to all keen students with any levels of experience.
        \newline
        
        \noindent
        Notably, while IRIS projects are not geared directly towards research groups, (depending on the proposed project) these projects may have the possibility of playing an indirect contribution, or even eventually tying to a Bachelors/Masters Thesis!
        \newline

        \noindent        
        Individuals and Groups (max. 4 students) are encouraged to apply!\footnote{Forms: Short 1-page IRIS application form; Detailed proposal using Chapters 1-5 of this document. (Sent to iris@lists.kit.edu)}
        
	\section{Tentative timeline}
	\begin{tabularx}{\textwidth} { 
			| >{\centering\arraybackslash\bf}X 
			| >{\raggedright\arraybackslash}X 
			| >{\raggedright\arraybackslash}X | }
		\hline
		Period / Date & \bf{Phase} & \bf{Comments} \\
		\hline
		  16 June - 1 July  & Publicity \& Application & -  \\
		\hline
		  1 July  & Proposal & Feedback about proposal  \\
		\hline
		1 June - 30 July  & Detailed proposal  & Details about resources  \\
		\hline
		30 July  & Project confirmation  & Approval by Professors  \\
		\hline
		1 August - 8 November  & Official IRIS Period  & Deliverables due by 8 November\\
		\hline
		
	\end{tabularx}

        \section{Further details}
        IRIS aims to be an outlet for students to start and execute their own projects, which inculcates multidisciplinary self-exploration and innovation under the Physics department of KIT. This is to be achieved through unique and novel short-term projects, proposed by students and identified through gaps with references.
        \newline

        \noindent
        Through direct involvement in self-directed research, students will hopefully gain and develop understanding for different aspects for projects, from research gap identification, to proposals, time and resource management.  They will also gain other soft skills which will be essential in future work (in and outside of academia), such as logistics management, critical thinking, setbacks management, verification testing, etc. Additionally, IRIS especially seeks to allow students to develop and gain confidence in their unique and highly individual skill-sets, while fulfilling milestones and completing tasks.
        \newline

        \noindent
        Groups are expected to self-direct all aspects of the project, notably, procurement details will need to be discussed as soon as possible. It is also important for students to convey the necessity for special access, such as chemical access, clean-room, or workshops. As well as for goods with long procurement lead-times.

       
	\section{Deliverables during IRIS Period}
	\subsection{Prototyping , Updates and Progress}
        At this stage, unique innovations are at the core of this initiative. These can come in many forms, robotics, written code, 3D printed samples, etc. It is strongly encouraged to have continual documentation (photos, videos, etc.) in some form of progress journal.\\
        
        \noindent 
        It is also recommended that prototyping should be an immediate priority, as well as continually making iterations and optimisations. For each prototype or draft, students should demonstrate continual improvements towards their proposed goal. The term "prototypes" here could also include related sub-components and crucial elements for the project.\\
        
        \subsection{Final presentation \& Demonstration day}
        There will be a final presentation where the teams demonstrate how their prototype fulfills milestones and goals they originally proposed in the beginning of their investigation (or present other results that refute or change the originally set goals). On the presentation day all teams and members need to be fully present. Teams must work their best to demonstrate a functional working prototype.
        \subsection{A2 Poster/ Article / Report}
        Generally, groups will have to work out a formal deliverable (In the form of a poster, article, report etc.) that captures the many factors that leads to the essence of the project. As well as elaborating on future continuation. \\
        The exact of this report is context-specific and will vary, since they will be tailored to each project. This final deliverable should contain all necessary information that outlines the project, and highlight the accomplishments that have been achieved.
       
%%%%%%%%%%%%%%%%%%%%%%%%%%%%%%%%%%%%%%%%%%%%%%%%%%%%%%%%%%%%%%%%%%%%%%%%%%%%%%%%%%%%%%%%%%%%%%%%%%%%%%%%%%%%%%%%%%%%%%%%%%%%%%%%%%%%%%%%%%%%%%%%%%%%%%%%%%%%%%%%%%%%%%%%%%%%%%%%%%%%%%%%%%%%%%%%%%%%%%%%%%%%%%%%%%%%%%%%%%%%%%%%%%%%%%%%%%%%%%%%%%%%%%%%%%%%%%%%%%%%%%%%%%%%%%%%%%%%%%%%%%%%%%%%%%%%%%%%%%%%%%%%%%%%%%%%%%%


      
 % ONLY FILL AND SIGN IN THE PAGES BELOW !! ONLY FILL AND SIGN IN THE PAGES BELOW !!  ONLY FILL AND SIGN IN THE PAGES BELOW !! ONLY FILL AND SIGN IN THE PAGES BELOW !! ONLY FILL AND SIGN IN THE PAGES BELOW  !!

 
 % DELETE THE PREWRITEN TEXT AND FILL IN THE FORM  DELETE THE PREWRITEN TEXT AND FILL IN THE FORM  DELETE THE PREWRITEN TEXT AND FILL IN THE FORM  DELETE THE PREWRITEN TEXT AND FILL IN THE FORM  

 %%%%%%%%%%%%%%%%%%%%%%%%%%%%%%%%%%%%%%%%%%%%%%%%%%%%%%%%%%%%%%%%%%%%%%%%%%%%%%%%%%%%%%%%%%%%%%%%%%%%%%%%%%%%%%%%%%%%%%%%%%%%%%%%%%%%%%%%%%%%%%%%%%%%%%%%%%%%%%%%%%%%%%%%%%%%%%%%%%%%%%%%%%%%%%%%%%%%%%%%%%%%%%%%%%%%%%%%%%%%%%%%%%%%%%%%%%%%%%%%%%%%%%%%%%%%%%%%%%%%%%%%%%%%%%%%%%%%%%%%%%%%%%%%%%%%%%%%%%%%%%%%%%%%%%%%%%%%%%%%%%%%%%%%%%
 
		
\chapter{Introduction of Proposed Project}

	\section{Proposed project \& Scope}
        The impetus for the project and its application would be motivated by the identified problem statement, with the constraints outlined by its declared scope.Teams should carefully state the scope that allows for a balance between adaptability and rigidity. This ensures that the project can be self-contained, and deemed completed within a set of constraints. 
        \newline

        \noindent
         It is important to remember that this is a short research and development time frame. There is a trade-off between how large the scope can be, to be able to fulfill the problem statement, and how narrow it must be, for the project to be completed. Constraints are necessary to avoid out-of-scope considerations by teams themselves, and also out-of-scope external feedback.

	\section{Current literature \& Gap }
 
        Teams need to provide current literature on the topic, with a minimum of 3 sources\footnote{Sources can include any sort of documentation, they need not be limited to scientific journals. Other examples may include, conference presentations, websites, online videos, patents, product spec-sheets, etc. Teams are also free to include any number of sources in this section, as long as they are able to identify a common problem statement and research gap relationship between all of them.} that are directly related to the proposed project in mind. These sources should be elaborated briefly and the common gap\footnote{The problem statement should be established, which may include examples such as: prohibitively high cost, lack of research, absence of invention, alternative methods, discontinued developments, reduction of barriers to entry, etc.} should be articulated clearly. The identified gap is the focal point where the proposed project would directly address, or it serves as an important part in a larger process. 
        \newline

        \noindent
        Teams may also include documentation on other attempts at realisation, with proper citations and inclusion of an annotated bibliography of the mentioned source. Teams are expected to practice their own discretion regarding sources used , and verify them according to their authenticity. Some possible sources where sources may come from are publicly available content, current investigations, or personal experiences. 
        
        \subsection{Source 1 - Placeholder Title}
        \lipsum[66]

        \subsection{Source 2 - Placeholder Title}
        \lipsum[75]

        \subsection{Source 3 - Placeholder Title}
        \lipsum[66]

%%%%%%%%%%%%%%%%%%%%%%%%%%%%%%%%%%%%%%%%%%%%%%%%%%%%%%%%%%%%%%%%%%%%%%%%%%%%%%%%%%%%%%%%%%%%%%%%%%%%%%%%%%%%%%%%%%%%%%%%%%%%%%%%%%%%%%%%%%%%%%%%%%%%%%%%%%%%%%%%%%%%%%%%%%%%%%%%%%%%%%%%%%%%%%%%%%%%%%%%%%%%%%%%%%%%%%%%%%%%%%%%%%%%%%%%%%%%%%%%%%%%%%%%%%%%%%%


\chapter{Foundation}

 	\section{Physical background concepts }
         There is a huge range of approaches for realizing the proposed project, generally they can be categorized into two different broad categories. Concepts important for cognition of the project will come from theory and principles of operation. Thereafter, methodology and processes will contribute steps that must be taken for realisation of these concepts.
    
        \subsection{Theory or Principles of operation }
        This section is akin to the section in a thesis that goes over background literature. It should be written with enough details to highlight critical parameters or modes of operation which is essential for the completion of prototypes (or parts). In this document, this subsection functions as a primer for important knowledge required for conceptualisation of the project.
        \newline

        \noindent
        Teams are free to add more subsections for separation of topics here. There is no need to elaborate on every possible related topic, only the essential ones within the constraints of the defined project scope. It is recommended that there should be a liberal amount of visuals, diagrams, or tables in this subsection as well. Avoid typing walls of text!


        \subsection{Methodology \& Process}
        Distinct from the previous section, this subsection will cover essential experimental and applied tools which are needed to make concepts into a reality. There is a need to provide some minimum of theoretical grounding or technical explanation regarding methodologies used, however this would not apply to all methodologies, only what is deemed essential or significant by students. Testing, analysis, and validation techniques that need to be use should also fall under this category. 
        \newline

        \noindent
        Some methodologies that teams may choose to use may not be open-source, but documentation on it's validity of use, in combination with elaboration on the use-case is sufficient. It should be evident to readers of this document to understand how contents in the next chapter come about. Which are resources such as procured goods, tools, laboratory access, chemical purchases, etc. (Should they apply).


%%%%%%%%%%%%%%%%%%%%%%%%%%%%%%%%%%%%%%%%%%%%%%%%%%%%%%%%%%%%%%%%%%%%%%%%%%%%%%%%%%%%%%%%%%%%%%%%%%%%%%%%%%%%%%%%%%%%%%%%%%%%%%%%%%%%%%%%%%%%%%%%%%%%%%%%%%%%%%%%%%%%%%%%%%%%%%%%%%%%%%%%%%%%%%%%%%%%%%%%%%%%%%%%%%%%%%%%%%%%%%%%%%%%%%%%%%%%%%%%%%%%%%%%%%%%%%%%%%%%%%%%%%%%%%%%%%%%%%%%%%%%%%%%%%%%%%%%%%%%%%%%%%%%%%%%%%%


\chapter{Resource management}


 	\section{Supporting physical resources }
        There are generally two basic aspects which can be mentioned here, goods and materials, as well as tools and equipment. It is important to highlight the essentials for both that are needed, as this will give indication on how much time is needed before materials and goods may arrive, as well as what laboratories or equipment that are permitted.
        \newline

        \noindent
        It is important to take note of lead-times, required licences, specialised laboratories, or other prerequisites. You may need to reach out and contact other parties, such as vendors for quotations, for example. (Or perhaps professors of labs, which may need to be accessed.)
        \newline

        \noindent
        It is encouraged to have a discussion well ahead of time to understand physical resources needed for realisation. Important questions such as "Where do we buy?", "How long ahead of time do we need to make the purchase?", or "How easy or difficult is it to obtain?" should come up.

        \subsection{Materials \& Goods}
        \lipsum[4]


        \subsection{Tools, Equipment, Labs}
        \lipsum[11]


	
	\section{Availability \& Task specialisation (Gantt Chart)}
        This section will primarily focus on details related to procurement and personnel. A Gantt chart will provide a succinct and comprehensive overview. Teams should include a final testing and streamlining phase in their Gantt chart, in the period before the final week of the IRIS period. This allows ease of viewing for both teams and organisers on who are available.
        \newline

        \noindent
        It is recommended to contain an approximate timeline, inclusive of dates before the IRIS period, with all tasks inside. Notably, procurement and certain applications would take place prior to the official start of IRIS. It is also good to state weeks or dates where team members are unavailable for IRIS.
        \newline

        \noindent % Put this in 4 ( Checkpoints)
        It is suggested to also be a good idea, to indicate team member's relevant experiences. It isn't essential to have exact experiences, as skills can always be learned\footnote{An example could be 3D modelling / 3D printing workflow knowledge, which may also include experience with troubleshooting experience. Skills such as this can be learned quickly by keen individuals, however pre-existing experience can provide an edge to certain projects that need it.}, it may also serve as an indication for feasibility of certain tasks.
	
%%%%%%%%%%%%%%%%%%%%%%%%%%%%%%%%%%%%%%%%%%%%%%%%%%%%%%%%%%%%%%%%%%%%%%%%%%%%%%%%%%%%%%%%%%%%%%%%%%%%%%%%%%%%%%%%%%%%%%%%%%%%%%%%%%%%%%%%%%%%%%%%%%%%%%%%%%%%%%%%%%%%%%%%%%%%%%%%%%%%%%%%%%%%%%%%%%%%%%%%%%%%%%%%%%%%%%%%%%%%%%%%%%%%%%%%%%%%%%%%%%%%%%%%%%%%%%%%%%%%%%%%%%%%%%%%%%%%%%%%%%%%%%%%%%%%%%%%%%%%%%%%%%%%%%%%%%%%%%%%%%%%%%%

\chapter{Milestones and Goals}


	\section{Checkpoints}
        It's important to identify crucial intermediate steps in the process of developing a prototype or draft. While these intermediate steps won't necessarily resemble the final iteration, anticipation and planning for important milestones can give important foresight which will give confidence in reaching the goal of the project within a reasonable amount of time. 
        \newline
        
        \noindent
        Setting these goals allow for anticipation of the workflow required to reaching a functional prototype. These checkpoints will allow for discussions prior to starting the project, and provide an avenue of progress-tracking.


	\section{First working draft / prototype \& Important components }   
        It's necessary to have foresight regarding a set of minimum requirements for an initial prototype (or just important components), that is within the project scope. Although subject to changes, a basic but cohesive prototype achieved early-on will demonstrate feasibility. Ideally, a first prototype / draft should be accomplished by the midpoint of the project. 
 
	
	\section{Contingency Plan/Points of Pivot}

    The project should maintain a degree of flexibility to accommodate unforeseen challenges, allowing for the possibility of pursuing secondary, adjacent, or tangential goals if necessary to ensure project success.


    %%%%%%%%%%%%%%%%%%%%%%%%%%%%%%%%%%%%%%%%%%%%%%%%%%%%%%%%%%%%%%%%%%%%%%%%%%%%%%%%%%%%%%%%%%%%%%%%%%%%%%%%%%%%%%%%%%%%%%%%%%%%%%%%%%%%%%%%%%%%%%%%%%%%%%%%%%%%%%%%%%%%%%%%%%%%%%%%%%%%%%%%%%%%%%%%%%%%%%%%%%%%%%%%%%%%%%%%%%%%%%%%%%%%%%%%%%%%%%%%%%%%%%%%%%%%%%%%%%%%%%%%%%%%%%%%%%%%%%%%%%%%%%%%%%%%%%%%%%%%%%%%%%%%%%%%%%%%%%%%%%%%%%%

\end{document}